% Complete Section 8.1, original PDF 107-108 / printed 97-98.
% block: eigenvalue-problems-heading page:107 type:section
\subsection{\BT{Eigenvalue problems}{特征值问题}}\label{sec:eigenvalue-problems}

% block: hodge-eigenvalue-problem page:107 type:ordinary-paragraph
\BT{Besides the Hodge Laplacian source problem, we are often interested in solving the associated eigenvalue problem}{除 Hodge Laplacian 的源问题外, 我们常常还希望求解相应的特征值问题}
\[
(\delta d+d\delta)u=\lambda u,
\]
\BT{together with boundary conditions, for $\lambda\in\mathbb R$, $u\neq0$. (Note that $\lambda=0$ is an eigenvalue if and only if there exist nonzero harmonic forms, and the harmonic forms form the eigenspace associated to the zero eigenvalue.)}{连同边界条件, 其中 $\lambda\in\mathbb R$, $u\neq0$. (注意, $\lambda=0$ 是特征值当且仅当存在非零调和形式, 并且调和形式构成零特征值对应的特征空间.)}

% block: hodge-eigenvalue-mixed-formulation page:107 type:ordinary-paragraph
\BT{Setting $\sigma=d^*u$, we obtain the mixed weak formulation of the eigenvalue problem: find $\lambda\in\mathbb R$, $\sigma\in H\Lambda^{k-1}$, $0\neq u\in H\Lambda^k$ such that}{令 $\sigma=d^*u$, 得到特征值问题的混合弱形式: 求 $\lambda\in\mathbb R$, $\sigma\in H\Lambda^{k-1}$, $0\neq u\in H\Lambda^k$, 使得}
\[
\begin{aligned}
\langle\sigma,\tau\rangle-\langle u,d\tau\rangle&=0,&\tau&\in H\Lambda^{k-1},\\
\langle d\sigma,v\rangle+\langle du,dv\rangle&=\lambda\langle u,v\rangle,&v&\in H\Lambda^k.
\end{aligned}
\]

% block: hodge-eigenvalue-finite-element-discretization pages:107,108 type:ordinary-paragraph
\BT{The finite element discretization is then clear. We replace the spaces $H\Lambda^{k-1}$ and $H\Lambda^k$ with a stable pair of spaces of finite element differential forms, $V_h^{k-1}$ and $V_h^k$, which we have derived for the source problem. This leads to a symmetric sparse generalized matrix eigenvalue problem, which can be solved by a variety of iterative eigensolvers, such as the Krylov–Schur method.}{此时有限元离散化就很清楚了. 将空间 $H\Lambda^{k-1}$ 和 $H\Lambda^k$ 替换为针对源问题已经导出的稳定有限元微分形式空间对 $V_h^{k-1}$ 和 $V_h^k$. 这就得到对称稀疏广义矩阵特征值问题, 可用 Krylov–Schur 方法等各种特征值迭代求解器求解.}

% block: eigenvalue-convergence-requirements page:108 type:ordinary-paragraph
\BT{The convergence theory for eigenvalue problems is quite subtle. A good convergence result must of course show that any of the (infinitely many) eigenvalues of the continuous problem is approximated arbitrarily closely by an eigenvalue of the discrete problem, if the mesh is sufficiently fine. It must also take into account multiplicity, so that an eigenvalue of the continuous problem of multiplicity $m$ is approximated by $m$ eigenvalues of the discrete problem (counting their multiplicity, which will usually be different from $m$). It must also rule out spurious eigenvalues, that is, show that the discrete eigenvalues do not accumulate, as the mesh is refined, at any number which is not an eigenvalue of the continuous problem. The approximation of eigenvectors has to be stated particularly carefully in the case of multiple eigenvalues.}{特征值问题的收敛理论相当精细. 一个好的收敛结果当然必须说明, 若网格充分细, 连续问题 (无穷多个) 特征值中的任意一个, 都能由离散问题的某个特征值任意精确地逼近. 还必须考虑重数, 使连续问题的一个 $m$ 重特征值, 由离散问题的 $m$ 个特征值逼近 (计入它们的重数, 该重数通常不同于 $m$). 还必须排除伪特征值, 即证明随着网格加密, 离散特征值不会在任何并非连续问题特征值的数处聚积. 对重特征值, 尤其必须谨慎陈述特征向量的逼近.}

% block: eigenvalue-convergence-analysis page:108 type:ordinary-paragraph
\BT{Fortunately, the theory of convergence of Galerkin approximations to eigenvalue problems has been carefully developed and it applies to the case of the Hodge Laplacian. The basic theory is due Babuška and Osborn \BookCite{16}. An excellent exposition is in \BookCite{18}. A brief discussion of its application to the Hodge Laplacian eigenvalue problem is given in [\BookNumber{11}, Section 8] and [\BookNumber{13}, Section 3.6] and so will not be repeated here. The main message is that the hypotheses of the Babuška–Osborn theory can be deduced from the improved error estimates of Theorem~\ref{thm:mixed-improved-errors}, and consequently whenever the hypotheses of that theorem---the compactness property of the Hilbert complex and $W$-bounded cochain projectors for the subcomplex of Galerkin spaces---are met, we obtain convergence of the eigenvalues and eigenvectors. The convergence rates for the eigenvectors are optimal (determined by the approximation properties of the spaces), and the convergence rates for the eigenvalues are double those of the eigenvectors.}{幸运的是, 特征值问题的 Galerkin 逼近收敛理论已得到细致发展, 并适用于 Hodge Laplacian. 基本理论由 Babuška 和 Osborn \BookCite{16} 建立. \BookCite{18} 给出了出色的阐述. [\BookNumber{11}, Section 8] 和 [\BookNumber{13}, Section 3.6] 简要讨论了它在 Hodge Laplacian 特征值问题中的应用, 这里不再重复. 主要结论是, Babuška–Osborn 理论的假设可由 Theorem~\ref{thm:mixed-improved-errors} 的改进误差估计推出. 因此, 只要该定理的假设, 即 Hilbert 复形的紧性性质以及 Galerkin 空间子复形的 $W$-有界上链投影得到满足, 就有特征值和特征向量的收敛性. 特征向量的收敛阶是最优的 (由空间的逼近性质决定), 而特征值的收敛阶是特征向量收敛阶的两倍.}

% block: hodge-and-maxwell-spectral-examples page:108 type:ordinary-paragraph
\BT{We have already seen an example of eigenvalue convergence for the Hodge Laplacian \eqref{eq:vector-laplacian-eigenproblem} in Table~\ref{tab:feec-vector-eigenvalues} of the introduction. By contrast, Table~\ref{tab:vector-lagrange-wrong-eigenvalues} showed that a naive method, not based on the mixed formulation, does not converge. In the introduction, we also considered the Maxwell eigenvalue problem $\operatorname{curl}\operatorname{curl}u=\lambda u$. A naive method based on Lagrange elements led to erroneous results, with the true spectrum either not recognizable from the computation or polluted by persistent spurious eigenvalues, as shown in Table~\ref{tab:maxwell-spectral-comparison}. Computations made with the weak formulation, but using the lowest order Whitney 1-forms, shown in the rows labeled FEEC in the table, converge nicely. This can be proven by relating the Maxwell eigenvalue problem to the Hodge Laplacian eigenvalue problem. The former arises as the $\mathcal B_1^*$ problem after a Hodge decomposition, and so the convergence analysis for the Hodge Laplacian eigenvalue problem carries over to the Maxwell problem. See [\BookNumber{13}, Section 3.6.1] for details.}{在引言的 Table~\ref{tab:feec-vector-eigenvalues} 中, 已经看到了 Hodge Laplacian \eqref{eq:vector-laplacian-eigenproblem} 特征值收敛的一个例子. 与此相反, Table~\ref{tab:vector-lagrange-wrong-eigenvalues} 表明, 一个不基于混合形式的朴素方法并不收敛. 引言中还考虑了 Maxwell 特征值问题 $\operatorname{curl}\operatorname{curl}u=\lambda u$. 基于 Lagrange 单元的朴素方法导致错误结果, 真正的谱或者无法从计算中辨认, 或者受到持续存在的伪特征值污染, 如 Table~\ref{tab:maxwell-spectral-comparison} 所示. 使用弱形式和最低阶 Whitney 1-形式进行的计算, 在表中标为 FEEC 的各行中给出, 收敛良好. 将 Maxwell 特征值问题与 Hodge Laplacian 特征值问题联系起来就可证明这一点. 前者在 Hodge 分解后作为 $\mathcal B_1^*$ 问题出现, 因而 Hodge Laplacian 特征值问题的收敛分析可推广到 Maxwell 问题. 细节见 [\BookNumber{13}, Section 3.6.1].}
